\documentclass[12pt,a4paper]{article}
\usepackage[a4paper,margin=2cm]{geometry}
\usepackage[T2A]{fontenc}
\usepackage[utf8]{inputenc}
\usepackage[russian]{babel}
\usepackage{amsmath,amssymb,mathtools}
\usepackage{array,booktabs,longtable}
\usepackage{xcolor}
\usepackage{hyperref}
\usepackage{tikz}

\hypersetup{
  colorlinks=true,
  linkcolor=blue!55!black,
  urlcolor=blue!55!black,
  pdftitle={Ревью домашней работы},
  pdfauthor={Лёвин Артём Александрович}
}
\setlength{\parindent}{0pt}
\setlength{\parskip}{0.55em}
\renewcommand{\arraystretch}{1.25}

% Краткое сообщение Екатерине:
% Екатерина, обратите внимание на раскрытие квадрата разности в задании 3,
% признаки делимости в задании 5 и разрядную запись трёхзначного числа в задании 8.
% В этих трёх местах важно доводить выбранную идею до проверяемого ответа.

\begin{document}

\hypersetup{pageanchor=false}
\begin{titlepage}
  \centering
  \vspace*{0.22\textheight}
  {\Huge\bfseries Ревью домашней работы\par}
  \vspace{2.2cm}
  {\Large Автор: Лёвин Артём Александрович\par}
  \vspace{1cm}
  {\large 14 сентября 2026 года\par}
  \vfill
  \begin{tikzpicture}[x=1cm,y=1cm]
    \draw[blue!45!black, line width=0.8pt]
      (-3.2,0) .. controls (-1.8,0.55) and (-0.8,-0.45) .. (0,0)
      .. controls (0.8,0.45) and (1.8,-0.55) .. (3.2,0);
  \end{tikzpicture}
  \vspace*{0.8cm}
\end{titlepage}
\hypersetup{pageanchor=true}
\pagenumbering{arabic}

\newpage
\section*{Сводная таблица}
\small
\setlength{\tabcolsep}{3pt}
\begin{tabular}{>{\centering\arraybackslash}p{0.07\linewidth}
                p{0.17\linewidth}
                p{0.25\linewidth}
                p{0.18\linewidth}
                p{0.18\linewidth}}
\toprule
\textbf{№} & \textbf{Тема} & \textbf{Суть ошибки} & \textbf{Ваш ответ} & \textbf{Правильный ответ} \\
\midrule
\hyperref[task:3]{3} & Квадрат разности, многочлены & Неверно раскрыт квадрат разности: ошибочны знак среднего члена и квадрат одночлена. & $16+30y-8y^2$ & $4y^2-18y+16$ \\
\addlinespace
\hyperref[task:5]{5} & Признаки делимости & Решение не доведено до математического отбора цифр по делимости на $2$ и $9$. & не указан & $98730$ \\
\addlinespace
\hyperref[task:8]{8} & Разрядная запись числа & Начат разбор записи числа, но не получено уравнение для первой и последней цифр и не найден минимум. & не указан & $701$ \\
\bottomrule
\end{tabular}
\normalsize

\newpage
\vspace*{\fill}
\begin{center}
  {\Huge\bfseries Разбор ошибок}
\end{center}
\vspace*{\fill}

\newpage
\phantomsection\label{task:3}
\section*{Задание 3}

\textbf{Текст задания.}
Представьте выражение
\[
(4-3y)^2-y(5y-6)
\]
в виде многочлена стандартного вида.

\textbf{Ваше решение.}
В записи квадрат разности был раскрыт так, что далее получилось
\[
16+24y-3y^2-5y^2+6y,
\]
а после приведения подобных слагаемых ---
\[
16+30y-8y^2.
\]

\textbf{Ваш ответ.}
\[
16+30y-8y^2.
\]

\textcolor{red}{\textbf{Ошибка:}} неверно раскрыт квадрат разности $(4-3y)^2$.

\textbf{Где именно возникает ошибка.}
Последний полностью корректный шаг --- исходная запись
\[
(4-3y)^2-y(5y-6).
\]
Первый неверный переход появляется при раскрытии $(4-3y)^2$. В полученной записи средний член имеет знак ``плюс'', а квадрат одночлена $3y$ фактически не превращается в $9y^2$.

\textbf{Почему это неверно.}
Для квадрата разности действует формула
\[
(a-b)^2=a^2-2ab+b^2.
\]
Знак среднего члена обязательно отрицательный, потому что в произведении
\[
(a-b)(a-b)
\]
два смешанных произведения равны $-ab$ и $-ab$, поэтому вместе они дают $-2ab$.

В данном выражении $a=4$, $b=3y$. Значит,
\[
(4-3y)^2=4^2-2\cdot4\cdot3y+(3y)^2.
\]
Здесь
\[
4^2=16,\qquad -2\cdot4\cdot3y=-24y,
\]
а
\[
(3y)^2=3^2y^2=9y^2.
\]
Поэтому правильное раскрытие квадрата разности имеет вид
\[
(4-3y)^2=16-24y+9y^2.
\]
Отдельно нужно раскрыть вторые скобки с учётом множителя $-y$:
\[
-y(5y-6)=-5y^2+6y.
\]

\textcolor{blue!60!black}{\textbf{Ключевая идея:}}
сначала без спешки применить формулу квадрата разности, отдельно раскрыть вторые скобки, и только после этого приводить подобные слагаемые.

\textbf{Исправление от последнего правильного шага.}
Начинаем с исходного выражения и исправляем именно первый ошибочный переход:
\begin{align*}
(4-3y)^2-y(5y-6)
&=16-24y+9y^2-5y^2+6y\\
&=(9y^2-5y^2)+(-24y+6y)+16\\
&=4y^2-18y+16.
\end{align*}

\textbf{Полное правильное решение.}
Используем формулу квадрата разности:
\[
(4-3y)^2=4^2-2\cdot4\cdot3y+(3y)^2
=16-24y+9y^2.
\]
Раскрываем вторые скобки:
\[
-y(5y-6)=-5y^2+6y.
\]
Подставляем оба результата:
\[
16-24y+9y^2-5y^2+6y.
\]
Приводим подобные слагаемые. Квадратные члены дают
\[
9y^2-5y^2=4y^2,
\]
линейные члены дают
\[
-24y+6y=-18y,
\]
постоянный член равен $16$. Получаем многочлен стандартного вида
\[
\boxed{4y^2-18y+16}.
\]

\textbf{Проверка.}
Подставим, например, $y=1$ в исходное выражение:
\[
(4-3)^2-1\cdot(5-6)=1-(-1)=2.
\]
Подставим $y=1$ в полученный многочлен:
\[
4-18+16=2.
\]
Значения совпали, что подтверждает преобразование.

\textcolor{teal!60!black}{\textbf{Итог:}}
правильный многочлен стандартного вида ---
\[
\boxed{4y^2-18y+16}.
\]

\textbf{Задача на закрепление.}
Представьте выражение
\[
(5-2x)^2-x(3x-4)
\]
в виде многочлена стандартного вида.

\newpage
\phantomsection\label{task:5}
\section*{Задание 5}

\textbf{Текст задания.}
Найдите наибольшее пятизначное число, которое делится на $18$ и у которого цифры расположены в порядке строгого убывания: каждая следующая цифра меньше предыдущей.

\textbf{Ваше решение.}
После номера задания записано: ``не помню''. Математический отбор цифр и проверка делимости не выполнены.

\textbf{Ваш ответ.}
Не указан.

\textcolor{red}{\textbf{Ошибка:}} решение не завершено; не применены одновременно признаки делимости на $2$ и на $9$.

\textbf{Где именно возникает ошибка.}
Последнего математического шага, который можно продолжить, в записи нет. Первый недостающий шаг --- перевести условие ``делится на $18$'' в два удобных требования к цифрам числа.

\textbf{Почему это важно.}
Так как
\[
18=2\cdot9
\]
и числа $2$ и $9$ взаимно просты, число делится на $18$ тогда и только тогда, когда оно делится одновременно на $2$ и на $9$.

Делимость на $2$ означает, что последняя цифра должна быть чётной.

Делимость на $9$ означает, что сумма всех цифр должна делиться на $9$.

Кроме того, требуется \emph{наибольшее} число. Поэтому цифры нужно подбирать слева направо, каждый раз стараясь взять максимально возможную цифру, но сохраняя возможность выполнить оба признака делимости.

\textcolor{blue!60!black}{\textbf{Ключевая идея:}}
максимизировать число слева направо и одновременно контролировать чётность последней цифры и сумму цифр, кратную $9$.

\textbf{Исправление.}
Самая большая возможная первая цифра --- $9$. Вторая при строгом убывании максимально может быть $8$. Третья максимально может быть $7$.

Ищем число вида
\[
987ab,
\]
где
\[
7>a>b\geq 0,
\]
а $b$ должно быть чётным.

Сумма первых трёх цифр равна
\[
9+8+7=24.
\]
Для делимости на $9$ нужно, чтобы
\[
24+a+b
\]
делилось на $9$. Так как при строгом убывании максимально возможны $a=6$ и $b=5$, вся сумма цифр не превосходит
\[
24+6+5=35.
\]
Значит, единственный возможный кратный $9$ результат между $24$ и $35$ --- это $27$. Поэтому
\[
a+b=3.
\]
Чтобы число было наибольшим, берём максимально возможную четвёртую цифру $a$. Из пар неотрицательных цифр с суммой $3$ и условием $a>b$ подходят $(3,0)$ и $(2,1)$. Пара $(3,0)$ даёт большее число, причём последняя цифра $0$ чётная. Получаем
\[
9>8>7>3>0.
\]

\textbf{Полное правильное решение.}
Полученный кандидат ---
\[
98730.
\]
Проверяем делимость на $2$: последняя цифра $0$, поэтому число чётное.

Проверяем делимость на $9$:
\[
9+8+7+3+0=27,
\]
а $27$ делится на $9$.

Следовательно, $98730$ делится на $18$.

Почему нельзя получить большее число? Цифры $9$, $8$, $7$ уже максимально возможны на первых трёх позициях. Для этих трёх цифр условие делимости на $9$ вынуждает две последние цифры иметь сумму $3$. Максимальная допустимая четвёртая цифра при этом равна $3$, а последняя --- $0$. Поэтому найденное число действительно наибольшее.

\textbf{Проверка.}
\[
98730:18=5485,
\]
то есть деление выполняется без остатка. Порядок цифр также строгий:
\[
9>8>7>3>0.
\]

\textcolor{teal!60!black}{\textbf{Итог:}}
наибольшее подходящее пятизначное число ---
\[
\boxed{98730}.
\]

\textbf{Задача на закрепление.}
Найдите наибольшее шестизначное число, которое делится на $18$ и у которого цифры расположены в порядке строгого убывания.

\newpage
\phantomsection\label{task:8}
\section*{Задание 8}

\textbf{Текст задания.}
Задумали трёхзначное число, последняя цифра которого не равна нулю. Из этого числа вычли число, записанное теми же цифрами в обратном порядке. В результате получили $594$. Найдите наименьшее трёхзначное число, обладающее таким свойством.

\textbf{Ваше решение.}
Записана идея вычитания числа вида $abc$ и числа с обратным порядком цифр $cba$, после чего выполнены отдельные пробные вычисления. Завершённого уравнения и окончательного ответа нет; далее записано: ``не помню''.

\textbf{Ваш ответ.}
Не указан.

\textcolor{red}{\textbf{Ошибка:}} правильная идея с обратной записью числа не переведена в разрядную формулу и не доведена до поиска минимального числа.

\textbf{Где именно возникает ошибка.}
Последний корректный шаг --- обозначение исходного числа через его цифры и понимание, что из него нужно вычесть число с обратным порядком цифр. Первый недостающий шаг --- записать оба числа через сотни, десятки и единицы.

\textbf{Почему это важно.}
Запись $abc$ в задаче про цифры нельзя использовать как обычное произведение букв. Если $a$, $b$, $c$ --- цифры трёхзначного числа, то само число равно
\[
100a+10b+c.
\]
Число с обратным порядком цифр равно
\[
100c+10b+a.
\]
Только после такой разрядной записи становится видно, что цифра десятков $b$ при вычитании сокращается.

\textcolor{blue!60!black}{\textbf{Ключевая идея:}}
записать число и его обратную запись через разряды; после сокращения десятков получить простое условие на первую и последнюю цифры.

\textbf{Исправление от последнего правильного шага.}
Пусть исходное число имеет цифры $a$, $b$, $c$. Тогда
\[
(100a+10b+c)-(100c+10b+a)=594.
\]
Сокращаем одинаковые десятки:
\begin{align*}
100a+c-100c-a&=594,\\
99a-99c&=594,\\
99(a-c)&=594.
\end{align*}
Делим обе части на $99$:
\[
a-c=6.
\]

\textbf{Полное правильное решение.}
Последняя цифра исходного числа не равна нулю, поэтому
\[
c\geq1.
\]
Из равенства
\[
a-c=6
\]
получаем
\[
a=c+6.
\]
Нужно найти \emph{наименьшее} трёхзначное число. В первую очередь нужно сделать минимальной цифру сотен $a$. Так как $c\geq1$, наименьшее возможное значение
\[
a=7,
\]
и тогда
\[
c=1.
\]

Цифра десятков $b$ в уравнении сократилась и на разность не влияет. Чтобы само число было наименьшим, берём наименьшую цифру десятков:
\[
b=0.
\]
Получаем число
\[
\boxed{701}.
\]

\textbf{Проверка.}
Число с обратным порядком цифр равно $107$. Вычитаем:
\[
701-107=594.
\]
Последняя цифра исходного числа равна $1$, то есть условие $c\ne0$ выполнено.

Также число действительно минимально: меньшая цифра сотен $a$ невозможна, потому что тогда при $a-c=6$ последняя цифра $c$ была бы равна нулю или стала бы недопустимой.

\textcolor{teal!60!black}{\textbf{Итог:}}
наименьшее подходящее трёхзначное число ---
\[
\boxed{701}.
\]

\textbf{Задача на закрепление.}
Задумали трёхзначное число, последняя цифра которого не равна нулю. Из этого числа вычли число, записанное теми же цифрами в обратном порядке, и получили $396$. Найдите наименьшее трёхзначное число, обладающее таким свойством.

\vfill
\begin{center}
\small Лёвин Артём Александрович эксклюзивно для Екатерины
\end{center}

\end{document}
